\chapter{Metodologia assistita da agenti AI}
\label{cap:metodologia-ai}

\intro{Il presente capitolo descrive in che modo è stata impiegata l'AI durante lo stage, incentrato sull'uso di un agente di codifica basato su un \gls{llm} come \emph{pair engineer} strutturato. Vengono illustrati gli ambiti di utilizzo, il workflow adottato, i benefici in termini di tempo e le mitigazioni adottate per i limiti noti.}\\

\section{Contesto e motivazione}

La pianificazione dello stage prevedeva una durata di 300 ore distribuite su 8 settimane,
a fronte di un ambito ampio: apprendimento di tecnologie nuove (OpenTelemetry \gls{sdk} .NET,
stack \emph{Grafana} LGTM, PromQL/TraceQL/LogQL, Serilog \gls{otlp}, \emph{Playwright}), realizzazione
di prototipi, strumentazione di più moduli del framework \emph{Sia.Web} e produzione della
relativa documentazione. Per colmare il gap di conoscenza in merito alle tecnologie e lo studio
della soluzione \emph{Sia.Web}, si è fatto ricorso all’utilizzo di agenti di intelligenza artificiale.
Questo approccio ha permesso di velocizzare le attività ripetitive e di replicare pattern già validati su nuovi moduli.
L'introduzione di agenti di codifica, come \emph{Claude Code}, ha reso più efficiente l'introduzione
di test automatici e sistemi di telemetria all’interno di un’applicazione di grandi dimensioni,
permettendo di creare, e soprattutto propagare, codesti processi di supporto, senza sottrarre eccessive risorse agli sviluppi principali.

\section{Ambiti di utilizzo}

L'agente è stato utilizzato in tre ambiti distinti, corrispondenti a quattro momenti differenti
dello stage.

\subsection{Studio delle tecnologie}

Nelle prime fasi, dedicate allo studio delle tecnologie, l'agente è stato impiegato
per sintetizzare la documentazione ufficiale delle varie tecnologie, e per confrontare le possibili
soluzioni da adottare. Tale attività ha permesso di ridurre in modo sensibile il tempo stimato di
\emph{onboarding}, concentrando lo studio sui soli punti
critici per il progetto invece che sulla documentazione integrale di ciascuna tecnologia.

\subsection{Prove di implementazione}

Durante la fase di progettazione, l'agente è stato usato per produrre rapidamente varianti di
\gls{poc} su scelte ancora da validare: che moduli testare, quale fosse
il carico computazionale introdotto, l'ambiente di hosting più conveniente (\emph{Docker} o servizi Windows).
Avere a disposizione codice compilabile in pochi minuti ha consentito di provare e
scartare alternative, invece di decidere a priori sulla base della sola documentazione.

\subsection{Propagazione della telemetria agli altri moduli}

L'ambito in cui l'uso dell'agente si è rivelato più produttivo è stato la propagazione
della strumentazione ai moduli del framework (\emph{Sia.Base}, \emph{Sia.Auth},
\emph{Sia.Grid}, \emph{Sia.Crud}, \ldots). Per evitare che la replicazione del pattern di
strumentazione su moduli diversi introducesse incoerenze (naming, cardinalità, attributi
di \emph{resource}), sono state definite \emph{Claude skill} dedicate:

\begin{itemize}
      \item \texttt{sia-telemetry-setup}: configurazione della pipeline OTel, \gls{samplerg},
            \emph{exporter}, \emph{appsettings};
      \item \texttt{sia-telemetry-traces}: \emph{span}, \emph{tag}, \emph{helper} \texttt{TrackQuery};
      \item \texttt{sia-telemetry-metrics}: metriche, \emph{naming}, cardinalità, pattern
            \texttt{RecordMetricError};
      \item \texttt{sia-telemetry-logs}: Serilog \gls{otlp}, correlazione \emph{log}-\emph{trace};
      \item \texttt{sia-telemetry-grafana}: \emph{dashboard} Grafana, query PromQL/TraceQL/LogQL,
            strumenti \gls{mcp}.
\end{itemize}
Ogni \emph{skill} è un contratto procedurale, non un semplice \emph{prompt}, che descrive
convenzioni, \emph{anti-pattern} e \emph{template} di codice per uno specifico aspetto. L'agente
carica la \emph{skill} appropriata quando riceve un task pertinente, garantendo che
moduli strumentati in momenti diversi rispettino le stesse convenzioni.

\subsection{Propagazione dei test automatici}

La stessa metodologia è stata replicata per l'introduzione e la propagazione dei test
automatici, sia lato \emph{backend} sia lato \emph{frontend}. In entrambi i casi è stato
definito un agente dedicato che codifica la struttura della \emph{suite}, lo stack
tecnologico di riferimento e le convenzioni di \emph{naming} e organizzazione dei file,
così che ogni nuovo test risulti coerente con quelli già esistenti.\\
Sul \emph{backend} l'agente copre le categorie di test unitari, \emph{smoke},
\emph{lifecycle} e \emph{summary}, ed è affiancato da due \emph{skill} specifiche per la creazione di
nuovi \emph{lifecycle test}, essendo quest'ultimo il pattern più ripetitivo della \emph{suite}.
Una \emph{skill} è orientata al modulo \gls{crud}, nella quale viene spiegato all'\gls{llm} a cosa serve quel test, vincoli specifici
dei test e dove raccogliere i dati necessari per la scrittura dei test; un'altra è specifica del modulo Grid e
ricopre le stesse funzioni della \emph{skill} precedente ma relative ai componenti griglia.
Sul \emph{frontend} l'agente è invece focalizzato sui test di usabilità
end-to-end basati su \emph{Playwright}, sull'organizzazione dei \emph{page object} e
sulla scrittura di scenari resilienti rispetto al rendering asincrono dei componenti.\\
Il beneficio osservato è analogo a quello registrato sulla telemetria: una volta
validato il pattern su un primo caso di riferimento, l'estensione a un nuovo componente
diventa una richiesta semantica all'agente, che in pochi minuti produce i file di
scenario e le dichiarazioni necessarie, senza che lo sviluppatore debba ricordare
convenzioni distribuite tra più punti del progetto.

\section{Workflow Agenti AI}
\label{sec:workflow-plan}

Per evitare i rischi tipici della generazione automatica di codice (allucinazioni, \emph{overfitting}
al primo caso visto), è stato adottato un \emph{workflow} strutturato in tre fasi.

\begin{enumerate}
      \item \textbf{Plan.} Per ogni task non banale, l'agente produce un file
            \texttt{plan.md} che descrive
            obiettivo, componenti coinvolti, modifiche pianificate file per file, assunzioni e
            domande aperte. Nessun codice viene prodotto in questa fase.
      \item \textbf{Confirm.} Lo sviluppatore revisiona il \texttt{plan.md}, può editarlo
            direttamente, integrarlo o rifiutarlo. L'agente si ferma \emph{esplicitamente} in
            attesa di approvazione.
      \item \textbf{Implement.} Solo dopo approvazione l'agente scrive un \texttt{task.md}, una
            \emph{checklist} operativa, nella quale vengono tracciati in ordine, i singoli interventi
            eseguiti dall'agente.
\end{enumerate}
Il \emph{workflow} produce un artefatto tracciabile per ogni intervento: \texttt{plan.md} e
\texttt{task.md} restano nel repository e fungono da registro decisionale.

\section{Confronto pianificazione vs esecuzione}

L'uso dell'agente ha prodotto un beneficio misurabile soprattutto nelle fasi a forte
componente ripetitiva (studio, \gls{poc}, propagazione del pattern di strumentazione), e un
beneficio meno marcato in quelle a forte componente decisionale (scelte architetturali,
analisi di cardinalità). Nel complesso, il tempo risparmiato nelle fasi assistite è stato
riinvestito nelle attività non automatizzabili: revisione semantica del codice prodotto,
validazione delle \emph{dashboard} contro Prometheus, affinamento iterativo delle
\emph{skill}.\\
Per dare un riferimento quantitativo a questo bilancio, lo studio delle tecnologie e lo
sviluppo dei \gls{poc} erano stati preventivati attorno alle 80 ore: con l'uso degli agenti,
il tempo effettivo si è ridotto a circa 30 ore. Una riduzione ancora più marcata si è
osservata nella propagazione dei test, ambito per il quale la pianificazione iniziale non
prevedeva nemmeno la copertura di tutti i moduli poi effettivamente strumentati: una
stima approssimativa indica che, senza il supporto degli \gls{llm}, il lavoro avrebbe
richiesto circa un mese, condensato invece in una settimana e mezza.

\section{Limiti e mitigazioni}

L'uso di un agente \gls{llm} comporta rischi noti che sono stati mitigati in modo sistematico.

\begin{description}
      \item[Allucinazioni su metriche/label.] L'agente tende a inventare nomi di metriche o
            \emph{label} plausibili ma inesistenti. \textbf{Mitigazione}: regola obbligatoria di
            validazione tramite \gls{mcp} Grafana, verificando l'esistenza di metriche e label in Grafana
            prima di chiudere ogni modifica delle dashboard.
      \item[Overfitting del prompt.] Le \emph{skill} iniziali, scritte a partire dal primo
            modulo strumentato (\texttt{Sia.Base}), contenevano riferimenti troppo specifici
            (nomi di metriche, \emph{tag}, stored procedure). \textbf{Mitigazione}: riscrittura
            delle \emph{skill} in termini generali, con placeholder \texttt{\{modulo\}} e
            riferimenti ai moduli canonici solo come esempi.
      \item[Perdita di contesto su sessioni lunghe.] La \emph{context window} dell'\gls{llm} è
            finita; su task lunghi, l'agente perde traccia di decisioni prese in precedenza.
            \textbf{Mitigazione}: memorizzazione persistente tramite \texttt{MEMORY.md} per
            progetto e per agente, \texttt{plan.md} e \texttt{task.md} come registro decisionale
            per ogni task.
      \item[Divergenza dei pattern tra moduli.] Sviluppatori diversi, o sessioni diverse, possono
            produrre variazioni del pattern di strumentazione. \textbf{Mitigazione}:
            centralizzazione dei pattern obbligatori nell'agente \texttt{sia-telemetry-developer}
            (es.\ counter errori controller via \texttt{IControllerDiagnostics} con Adapter, uso
            obbligatorio della \emph{template variable} \texttt{bucket} nelle dashboard).
\end{description}
Per superare gran parte di queste limitazioni e arricchire i processi di ragionamento degli agenti AI
con contenuti e contesto aggiuntivi, è stato adottato il protocollo \gls{mcp}.
Considerato inoltre il crescente impatto di questo strumento nell’ecosistema \gls{llm},
si è ritenuto opportuno dedicargli una sezione specifica di approfondimento.


% =============================================================================
\section{Model Context Protocol}
\label{sec:mcp-methodology}
% =============================================================================
Il \gls{mcp} è il protocollo, introdotto nel 2024, che consente a un agente basato su modelli
linguistici di interagire con servizi esterni seguendo uno standard, ottenendo dati dinamici e
contestuali.

\subsection{Il problema che MCP risolve}

Un modello linguistico di grandi dimensioni \gls{llm} è, nella
sua forma base, un sistema chiuso: conosce ciò che è stato scritto fino alla sua data di
addestramento e risponde in linguaggio naturale, ma non può aprire un file, interrogare un
database, chiamare un'\gls{api} o agire nel mondo esterno. Questa limitazione non dipende
dall'intelligenza del modello, dipende dall'assenza di un canale strutturato attraverso
cui il modello possa ricevere dati aggiornati e restituire azioni concrete.\\
Prima dell'avvento di \gls{mcp}, ogni integrazione tra un \gls{llm} e un sistema esterno richiedeva
un connettore ad hoc: codice scritto una volta per quella specifica combinazione di
modello e strumento, da riscrivere ogni volta che cambiava uno dei due. Con la
proliferazione dei modelli e degli strumenti, il costo di manutenzione di questa matrice
di integrazioni punto-a-punto è diventato rapidamente insostenibile.\\
Il \gls{mcp}, introdotto da Anthropic nel novembre 2024 e
oggi standard open-source governato dalla Linux Foundation sotto l'egida dell'Agentic AI
Foundation, risolve il problema della standardizzazione definendo un'interfaccia universale che qualsiasi modello e
qualsiasi strumento possono implementare una volta sola per comunicare con l'intero
ecosistema.

\subsection{Dettagli del protocollo}
\label{sec:mcp-dettagli}

\gls{mcp} è un protocollo client-server basato su \emph{JSON-RPC 2.0} con
connessioni \emph{stateful}, articolato su tre ruoli: l'\textbf{host}
(l'applicazione che ospita l'\gls{llm}, responsabile del consenso dell'utente
e del controllo degli accessi), il \textbf{client} (il componente che gestisce
la connessione verso un server specifico) e il \textbf{server \gls{mcp}} (il
servizio, locale o remoto, che espone dati e funzionalità all'\gls{llm}).
Prima di qualsiasi scambio applicativo, client e server negoziano le proprie
capacità in una fase di \emph{handshake}, garantendo la retrocompatibilità
tra versioni diverse del protocollo.\\
Sul piano funzionale, il server può esporre tre tipi di primitive: le
\textbf{Resources} (dati consultabili, identificati da URI), i \textbf{Tools}
(funzioni invocabili dal modello per agire nel mondo, che richiedono consenso
esplicito dell'utente) e i \textbf{Prompts} (template predefiniti per
workflow ricorrenti). Simmetricamente, il client può offrire al server tre
capacità: il \textbf{Sampling} (che delega all'utente il controllo su quali
prompt vengono inviati al modello), i \textbf{Roots} (i confini del filesystem
o degli URI entro cui il server può operare) e l'\textbf{Elicitation} (la
possibilità per il server di richiedere informazioni aggiuntive in modo
strutturato).\\
Data l'ampiezza dei permessi che un server \gls{mcp} può richiedere, la
specifica affronta la sicurezza attraverso quattro principi: consenso
esplicito dell'utente prima di ogni invocazione di tool (con cautela verso le
descrizioni provenienti da server non verificati, vettore noto di
\emph{prompt injection}); privacy dei dati, che non devono essere trasmessi
ai server senza consenso; controllo sul sampling, con visibilità
dell'utente sui prompt inviati al modello; e separazione dei contesti, per
limitare l'esposizione di informazioni confidenziali. Ricerche indipendenti
hanno comunque evidenziato vulnerabilità concrete a \emph{tool poisoning} e
\emph{prompt injection} in implementazioni permissive, da cui l'indicazione
di adottare meccanismi di autorizzazione granulari.

\subsection{Applicazione nel progetto}

Nel contesto del progetto \emph{Sia.Core}, \gls{mcp} è stato impiegato per collegare l'agente
di sviluppo direttamente a Grafana, a \emph{Playwright} e a \emph{mssql}, eliminando il
passaggio manuale di schemi, nomi di metriche e strutture dati tra l'utente e l'agente.\\
L'impatto più evidente si è registrato nella generazione delle dashboard: prima
dell'integrazione, l'agente produceva widget plausibili ma spesso errati, metriche
inesistenti, attributi di filtraggio non presenti nelle serie temporali, tipi di grafico
inadatti alla natura del dato. Connettere l'agente a Grafana tramite \gls{mcp} ha ridotto
sensibilmente questi errori, perché il modello poteva interrogare direttamente l'istanza
e ricevere come contesto i nomi reali delle metriche, le label disponibili e la struttura
delle trace esistenti in Tempo.\\
Nonostante il miglioramento, si è resa necessaria una fase di \emph{prompt engineering}
mirata a consolidare le linee guida operative dell'agente. In particolare, sono state
introdotte direttive esplicite su tre fronti: quali tipi di grafico associare a ciascuna
categoria di metrica (es.\ \emph{time series} per latenze e rate, \emph{stat} per
contatori cumulativi, \emph{histogram} per distribuzioni di percentili); quali variabili
di template definire a livello di dashboard per abilitare il filtraggio per
\texttt{service\_name}, \texttt{sia\_repository} e \texttt{sia\_operation}; come
organizzare la griglia dei widget per mantenere leggibilità e coerenza visiva tra moduli
diversi.\\
Un beneficio analogo si è ottenuto con l'integrazione di \emph{Playwright} e
\emph{mssql}: in entrambi i casi, l'accesso diretto agli strumenti tramite \gls{mcp} ha
permesso all'agente di operare su dati reali anziché su assunzioni, riducendo il numero
di iterazioni necessarie per ottenere un risultato corretto.
L'accesso al database ha permesso all'\gls{llm} di apportare modifiche direttamente a tabelle,
viste e \emph{stored procedure}, oltre a recuperare dati utili per il \emph{debugging} o per
ricostruire la struttura delle entità coinvolte, condizione necessaria per scrivere query
corrette senza doverla descrivere manualmente a ogni richiesta. Il collegamento di
\emph{Playwright} tramite \gls{mcp} ha reso possibile all'\gls{llm} non solo di eseguire
direttamente i test E2E, ma anche di aprire schede del browser per validare visivamente le
modifiche apportate e ispezionare bug grafici, comprendendone la causa e intervenendo di
conseguenza.

\section{Considerazioni conclusive}

L’utilizzo di un agente AI come supporto allo sviluppo, insieme all’adozione del protocollo \gls{mcpg},
ha contribuito in modo significativo al miglioramento della qualità del prodotto,
alla riduzione dei tempi di sviluppo e al superamento delle iniziali lacune tecniche.\\
Tuttavia, il ruolo dell’agente AI non deve essere interpretato come sostitutivo di quello dello sviluppatore.
Il lavoro eseguito dall'agente, che fosse codice prodotto, scelte architetturali
o propagazione degli sviluppi, ha spesso richiesto interventi di indirizzamento, correzione e validazione
da parte dello sviluppatore risultando quindi un processo collaborativo piuttosto che autonomo.
